% Generated by Listas/make_statements.py from pratica-lecture-08.md, item 16. Do not edit.
\textbf{Source.} Tadelis (2013), \emph{Game Theory: An Introduction}, Exercise 16.8 ``Beer or Quiche?'', p.~338 {[}PDF p.~357{]}. Reproduced verbatim. Answer: none in the book.



\begin{quote}

\textbf{Beer or Quiche?} Cho and Kreps (1987) introduced what is now a famous two-player signaling game. First, Nature selects player 1, who knows his type, to be either a wimp (\(W\)), with probability \(p = 0.1\), or surly (\(S\)), with probability \(1 - p = 0.9\). Player 1 then chooses what to have for breakfast: beer (\(B\)) or quiche (\(Q\)). A surly type prefers beer while a wimp prefers quiche. Player 1's preferred breakfast gives him a payoff of 1 while his less-preferred choice gives him 0. After breakfast player 2 observes what player 1 ate but does not know whether he is a wimp or surly. Player 2 then chooses whether to duel (\(D\)) with player 1 or not to duel (\(N\)). Player 1, regardless of his type, prefers no duel, yielding him an extra payoff of 2, to a duel, which gives him 0. (For example, if player 1 eats his preferred breakfast and avoids a duel then his final payoff is 3, while if he eats his preferred breakfast and is forced into a duel then his final payoff is 1.) Player 2, however, prefers to duel if and only if player 1 is a wimp. If player 1 is surly then player 2's payoff is 0 from \(D\) and 1 from \(N\). If player 1 is a wimp then player 2's payoff is 2 from \(D\) and 1 from \(N\).



\begin{enumerate}

\def\labelenumi{\alph{enumi}.}

\item

  Draw the extensive form of this game.

\item

  What are the Bayesian Nash equilibria of this game?

\item

  What are the perfect Bayesian equilibria?

\item

  Of the equilibria you found in (c), which fail the intuitive criterion?

\end{enumerate}

\end{quote}



Note (errata, \texttt{.agent/notes/errata-game-theory-books.md}): Tadelis's Def. 16.1 (the intuitive criterion) prints \(v_2\) in condition 2 where \(v_1\) belongs; apply the corrected version (\texttt{Notas/L08/L08-08-intuitive-criterion.md}).

